% 原书分章PDF第1—4页：章首与第27.1节。
\chapter{敏感性分析}
\label{ch:sensitivity-analysis}
\chapterauthors{Iván Díaz、Alexander R. Luedtke 与 Mark J. van der Laan}
\chapteraffiliations{%
I. Díaz（通讯作者）：美国康奈尔大学威尔康奈尔医学院医疗保健政策与研究系生物统计学及流行病学部，402 East 67th Street, New York, NY 10065, USA；电子邮箱：\href{mailto:ild2005@med.cornell.edu}{ild2005@med.cornell.edu}。\\
A. R. Luedtke：美国弗雷德·哈钦森癌症研究中心疫苗与传染病部，1100 Fairview Ave, Seattle, WA 98109, USA；电子邮箱：\href{mailto:aluedtke@fredhutch.org}{aluedtke@fredhutch.org}。\\
M. J. van der Laan：美国加利福尼亚大学伯克利分校生物统计学部及统计学系，101 Haviland Hall, \#7358, Berkeley, CA 94720, USA；电子邮箱：\href{mailto:laan@berkeley.edu}{laan@berkeley.edu}。\\
本章 DOI：\url{https://doi.org/10.1007/978-3-319-65304-4_27}。}

因果推断问题往往通过研究一列反事实变量分布的参数来处理；这些变量表示在施加某种干预的假想世界中的结局。因果推断的根本困难在于，对于一个给定个体，我们只能观测其中一个反事实结局：实际所分配处理水平下的结局。因此，要从观测数据的分布识别缺失反事实结局的分布，就必须作出某些不可检验的假设。其中一个常见假设是：处理机制不依赖于与结局存在因果联系的未测量因素。这个假设通常被称为处理分配的不可忽略性（nonignorability）（Rubin，1976），或（序贯）随机化假设（van der Laan 和 Robins，2003）。

\begin{sourcebox}
如果依据当前科学知识，可识别性假设无法得到合理支持，那么就不能从观测数据估计反事实分布。在这种情况下，一种有用的方法是进行敏感性分析（sensitivity analysis），帮助研究者了解：当假设受到不同程度的违背时，分析结果会如何改变。
\end{sourcebox}

关于因果推断的敏感性分析，已有大量文献。其中一种最常用的方法，由 Rotnitzky 等（1998）、Scharfstein 等（1999）、Robins 等（1999）、Robins（1999）、Rotnitzky 等（2001）及 Scharfstein 和 Robins（2002）提出。他们的方法假定：给定反事实结局后，处理机制的条件分布属于某个参数模型，该模型由一个不可识别的参数 $\alpha$ 索引。其敏感性分析是在假定的不同 $\alpha$ 水平下估计处理效应。随后，利用领域科学知识判断每个 $\alpha$ 值是否合理，并将合理 $\alpha$ 值对应的因果参数估计作为研究结论。

在 Díaz 和 van der Laan（2013b）以及 Luedtke 等（2015b）中，我们讨论了上述方法，并指出其一些局限。上一段所述方法，依赖于研究者能够正确地对反事实结局与处理分配机制之间的关系设定参数模型。在这一步使用参数模型，至少会对结论的有效性造成三种威胁。第一，人们早已认识到，参数模型不足以描述中高维数据中出现的复杂关系（例如见 Starmans，2011）。第二，即使只将参数模型用作一种近似，这种近似是否适当仍无法检验，因为反事实变量不可观测。第三个问题是所有参数模型共同面对的：在模型设定错误时，$\alpha$ 不再具有可理解的解释。因此，敏感性分析的支柱随之坍塌，因为我们无法根据辅助科学知识判断某些 $\alpha$ 值是否合理。

\begin{sourcebox}
本章提出一种能够克服上述问题的敏感性分析方法。这一方法与部分识别（partial identification）文献密切相关，可概括如下。我们将偏差分为两个主要分量：识别偏差（identification bias），定义为观测数据参数与因果参数之差；统计偏差（statistical bias），定义为估计量在重复抽样下的期望与观测数据参数之差。我们将讨论如何分别处理这两种偏差来源。对于识别偏差，我们建议使用因果效应的非参数界，这些界通常由某个可解释的条件 $\delta$ 索引。对于统计偏差，我们建议采用靶向学习（targeted learning）框架加以减小。
\end{sourcebox}

\section{问题}

本章始终使用以下两个实例说明相关概念。

\begin{bookresult}{例}{27.1}
Díaz 和 van der Laan（2013b）给出了一项展示恰加斯病（Chagas disease）治疗有效性的研究。由于该病潜伏期长达30年，必然会产生信息性退出（informative dropout），此前人们认为无法在这一应用中得出因果结论。如此长的潜伏期也使随机化研究的费用高得难以承受，因此研究者只能利用观察性数据评价治疗的疗效。该问题最初的分析过程如下。对恰加斯病文献的全面回顾，找到了19项研究，合计约520名患者。对这些研究进行标准荟萃分析（meta-analysis）被认为并不适当，因为：（a）许多研究没有对照组；（b）所有文章中关于处理分配的信息都很有限；（c）没有研究报告基线或混杂因素信息，例如基线健康状况、病媒暴露等；（d）许多研究中都有大量失访和退出。就处理分配而言，上述各点意味着这些研究不能被视为随机化研究，而且可能存在相当程度的未测量混杂，例如根据未测量的基线感染状况分配治疗。此外，患者可能因健康状况恶化或好转而退出研究，因此失访与退出通常和未观测结局有关。
\end{bookresult}

\begin{bookresult}{例}{27.2}
Luedtke 等（2015b）分析了西部协作组研究（Western Collaborative Group Study，WCGS）的数据（Rosenman 等，1964，1975）。WCGS 是一项前瞻性队列研究，旨在了解二分类人格类型（A型或B型）对八年半内冠心病（coronary heart disease，CHD）的影响。这些数据可以通过 R 的 \texttt{epitools} 包公开获取（Aragon，2012）。作者关注吸烟状态（是或否）对冠心病的影响。这个实例特别适合用于检验敏感性分析方法，因为文献已经充分确立了吸烟与冠心病之间的因果联系。Ding 和 VanderWeele（2015）在利用一份探究吸烟对肺癌影响的历史数据集评价其敏感性分析程序时，也作了类似选择。用来作为基线吸烟状态与冠心病之间关系的潜在混杂因素的协变量包括：年龄、身高、体重，以及A型人格的指示变量。我们考虑的结局是基线后 $8\frac12$ 年内是否发生冠心病。
\end{bookresult}

现在正式定义推断问题。为简单起见，我们聚焦于结局可能缺失的横断面研究。令 $A$ 表示二分类处理变量，$Y$ 表示连续或二分类结局，且仅在 $\Delta=1$ 时被观测；令 $W$ 表示已观测的处理前协变量向量。令 $O=(W,A,\Delta,\Delta Y)$ 表示分布为 $P_0$ 的随机变量，并令 $O_1,\ldots,O_n$ 表示 $O$ 的 $n$ 个独立同分布观测构成的样本。假定 $P_0\in\mathcal M$，其中非参数模型 $\mathcal M$ 定义为 $O$ 上相对于某个控制测度 $\nu$ 的所有连续密度。例27.1和27.2是这一设定的特例。在例27.1中，没有观测到协变量 $W$。在例27.2中，对 $(A,W)$ 的抽取以概率1满足 $P_0(\Delta=1\mid A=a,W=w)=1$，因此没有结局缺失。

在上述例27.1和27.2中，我们关心因果参数的估计；这些参数依据反事实结局 $Y_a:a\in\{0,1\}$ 的分布定义。这些结局表示在一个满足 $P_0(A=a,\Delta=1)=1$ 的假想世界中所观测到的结局。现在考察两个实例中的参数和可识别性假设。

\begin{bookresult}{例}{27.1（续）}
感兴趣的因果参数是已处理者中的处理效应，定义为 $\psi_{\mathrm{causal}}=E(Y_1-Y_0\mid A=1)$。在以下假设下：
\begin{enumerate}
\item $Y=\Delta(AY_1+(1-A)Y_0)$（一致性）；
\item $(Y_0,Y_1)\indep(A,\Delta)$（随机化与完全随机缺失）；
\item 对 $a\in\{0,1\}$，$P_0(A=a,\Delta=1)>0$（正值性），
\end{enumerate}
可以将 $\psi_{\mathrm{causal}}$ 识别为
\begin{equation}
 E(Y_1-Y_0\mid A=1)=E(Y\mid\Delta=1,A=1)-E(Y\mid\Delta=1,A=0),
 \tag{27.1}
\end{equation}
其中右侧参数可由观测数据估计。随机化与随机缺失意味着，处理机制和缺失机制独立于反事实结局。由于处理并非随机分配，而且认为退出研究与疾病严重程度有关，假设（2）很可能受到违背。
\end{bookresult}

\begin{bookresult}{例}{27.2（续）}
本问题的因果参数是平均处理效应，写为 $\psi_{\mathrm{causal}}=E(Y_1-Y_0)$。所需假设是根据观察性数据进行因果推断时的标准假设，即：
\begin{enumerate}
\item $Y=AY_1+(1-A)Y_0$（一致性）；
\item $(Y_0,Y_1)\indep A\mid W$（随机化）；
\item 对 $a\in\{0,1\}$，$P_0(A=a\mid W=w)>0$，关于 $W$ 的抽取以概率1成立（正值性）。
\end{enumerate}
在这些假设下，可写出 $E(Y_a)=E\{E(Y\mid A=a,W)\}$，其中右侧参数只依赖于观测数据 $O=(W,A,Y)$ 的分布。
\end{bookresult}
